\documentclass[11pt,a4paper]{article}

% ---- Packages ----
\usepackage[utf8]{inputenc}
\usepackage[T1]{fontenc}
\usepackage{times}
\usepackage{amsmath}
\usepackage{amssymb}
\usepackage{graphicx}
\usepackage{hyperref}
\usepackage{url}
\usepackage{booktabs}
\usepackage{geometry}
\usepackage{natbib}
\usepackage{microtype}
\usepackage{xcolor}
\usepackage{algorithm}
\usepackage{algpseudocode}
\usepackage{enumitem}
\usepackage{multirow}
\usepackage{array}

\geometry{
  left=2.5cm,
  right=2.5cm,
  top=2.5cm,
  bottom=2.5cm
}

\hypersetup{
  colorlinks=true,
  linkcolor=blue,
  citecolor=blue,
  urlcolor=blue
}

% ---- Title ----
\title{%
  \textbf{From Plausibility to Permissibility:\\
  A Layered Neuro-Symbolic Runtime for Reliable\\
  and Auditable AI Interpretation}
}

\author{%
  ANIMA \\
  ANIMA Research \\
}

\date{}

\begin{document}

\maketitle

% ---- Abstract ----
\begin{abstract}
Recent generative AI systems can produce highly plausible interpretations of linguistic and operational inputs, yet plausibility alone does not guarantee execution admissibility. In real-world settings, a fluent or high-confidence output may still be semantically ill-typed, structurally inconsistent, policy-violating, or insufficiently justified for action. This paper presents a layered neuro-symbolic interpretation runtime that explicitly separates candidate generation from execution eligibility.

The proposed framework combines candidate proposal, typed semantic role filtering, well-formedness constraints, policy arbitration, action gating, and audit tracing. Candidate interpretations are generated by a proposal layer, filtered by semantic and structural admissibility, and assigned one of three operational outcomes: \texttt{ACCEPT}, \texttt{REVIEW}, or \texttt{REJECT}. Hard constraints are applied prior to score-based ranking, preventing high-confidence but invalid candidates from being selected. Guarded override and conservative fallback mechanisms further reduce unsafe execution under uncertainty.

On an internal calibration set (20 cases across clear-accept, clear-reject, review-boundary, and override-mixed families), the runtime achieved 20/20 decision-match accuracy with explicit trace outputs preserved for every decision. These results establish internal consistency of the architecture under controlled operational conditions. The central implication is that reliable AI action requires post-generative control layers that transform plausible outputs into semantically valid, policy-aligned, and auditable decisions.

\vspace{6pt}
\noindent\textbf{Subtitle:} Neural Proposal, Symbolic Guarantee: Integrating Parser Policy, Typed Semantic Frames, and Policy Gating for ANIMA-Class Interpretation Runtimes

\vspace{6pt}
\noindent\textbf{Keywords:} neuro-symbolic runtime, interpretation admissibility, policy gating, typed semantic frames, auditability, plausibility vs.\ permissibility
\end{abstract}

% ==================================================================
% 1. INTRODUCTION
% ==================================================================
\section{Introduction}

Artificial intelligence systems are increasingly expected not only to interpret language, but also to authorize actions, invoke tools, and participate in operational workflows. This shift changes the deployment criterion. In many practical systems, the core question is no longer whether a model can produce a plausible interpretation of input, but whether that interpretation is permissible for execution under semantic, authority, and policy constraints.

A model output may be fluent, persuasive, and statistically likely while still failing in deployment-critical ways:
\begin{itemize}[leftmargin=*, itemsep=2pt]
  \item assigning impossible roles,
  \item collapsing distinct meanings,
  \item violating safety rules,
  \item exceeding authority boundaries,
  \item lacking auditable decision grounds.
\end{itemize}

This is especially acute in language-mediated systems where ambiguity is structural rather than exceptional. Natural language routinely contains attachment ambiguity, scope ambiguity, underspecified authority, and context-sensitive intent. Modern generative models are strong at proposing candidate interpretations under these conditions, but proposal quality and execution admissibility are different problems.

This work addresses that boundary directly. We propose a layered runtime in which a generative layer proposes candidate meanings, but explicit control layers determine whether a candidate is semantically valid, structurally well formed, policy-aligned, and auditable.

The central claim is:

\begin{quote}
\emph{Plausibility and permissibility are distinct layers of competence.}
\end{quote}

If this claim is correct, scaling generative capability alone is insufficient for trustworthy deployment. Systems require explicit post-generative control layers that convert candidate interpretations into admissible decisions.

\paragraph{Contributions.} This paper makes four contributions:
\begin{enumerate}[leftmargin=*, itemsep=2pt]
  \item It formalizes the distinction between plausible interpretation and permissible execution as a systems-level problem.
  \item It introduces a layered neuro-symbolic runtime integrating typed semantics, structural constraints, policy gating, and auditability.
  \item It provides an evaluation protocol that includes unsafe acceptance, review-boundary behavior, and trace completeness, not only selection accuracy.
  \item It reports internal calibration evidence showing stable and traceable decision behavior under policy-sensitive conditions.
\end{enumerate}

% ==================================================================
% 2. RELATED WORK
% ==================================================================
\section{Related Work}

\subsection{Neural Language Systems and Reliability}
Large language models have shown broad competence in instruction following, QA, summarization, planning, and code support. However, high output quality does not automatically guarantee semantic admissibility, safety compliance, or operational authority correctness when outputs are directly coupled to actions.

Calibration and uncertainty methods improve behavior but do not fully solve execution admissibility: a high-confidence output can still be operationally impermissible.

\subsection{Formal Semantics and Typing}
Formal semantics, typed roles, and frame-based representations provide mechanisms to separate grammatical possibility from semantic admissibility. Typed constraints are particularly effective for filtering role assignments that are linguistically plausible but semantically invalid.

\subsection{Neuro-Symbolic Integration}
Neuro-symbolic systems commonly combine neural proposal strengths with symbolic control strengths. Prior work emphasizes reasoning and compositionality benefits. This paper focuses on a deployment runtime objective: reliable and auditable action selection in language-mediated operational settings.

\subsection{Trustworthy AI and Auditability}
Trustworthy AI literature emphasizes safety, transparency, and governance. In operational interpretation systems, auditability requires explicit runtime decision structure: what was proposed, what was rejected, why, and under which rules.

% ==================================================================
% 3. PROBLEM FORMULATION
% ==================================================================
\section{Problem Formulation}

Given an input $x$, a proposal model generates candidate interpretations:
\[
I = \{i_1, i_2, \dots, i_n\}.
\]

A direct-selection baseline typically ranks candidates by plausibility and selects the top candidate. We instead define executable interpretation through four properties:
\begin{enumerate}[leftmargin=*, itemsep=2pt]
  \item \textbf{Plausibility}: statistical/linguistic reasonableness.
  \item \textbf{Semantic validity}: role and type compatibility.
  \item \textbf{Operational permissibility}: policy, authority, and safety compliance.
  \item \textbf{Auditability}: explicit traceability of decision and rejection path.
\end{enumerate}

Plausibility may guide proposal, but admissible execution requires semantic validity, permissibility, and auditability.

% ==================================================================
% 4. LAYERED RUNTIME ARCHITECTURE
% ==================================================================
\section{Layered Runtime Architecture}

\subsection{Overview}
The runtime consists of five sequential layers:
\begin{enumerate}[leftmargin=*, itemsep=2pt]
  \item \textbf{Proposal Layer}: generate interpretation candidates.
  \item \textbf{Typed Semantic Layer}: filter by role-type admissibility.
  \item \textbf{Well-Formedness Layer}: enforce structural and semantic coherence.
  \item \textbf{Policy Layer}: arbitrate with soft/hard rules into \texttt{ACCEPT}/\texttt{REVIEW}/\texttt{REJECT}.
  \item \textbf{Audit Layer}: preserve full trace for inspection and calibration.
\end{enumerate}

The governing principle is:

\begin{quote}
\emph{Generation proposes; structured control disposes.}
\end{quote}

\subsection{Typed Semantic Frames}
Candidates are mapped into typed event-role frames. Example templates:
\begin{itemize}[leftmargin=*, itemsep=2pt]
  \item \texttt{SEE(event)} with \texttt{agent: Human}, \texttt{patient: Entity}, \texttt{instrument: OpticalDevice?}
  \item \texttt{EXPLAIN(event)} with \texttt{agent: Human | System}, \texttt{patient: Claim | Anomaly}
  \item \texttt{OPEN(event)} with \texttt{agent: Human | System}, \texttt{patient: SecureObject}, \texttt{instrument: KeyLike?}
\end{itemize}

Type-incompatible role assignment is immediately rejected.

\subsection{Well-Formedness Constraints}
We define four constraints:
\begin{itemize}[leftmargin=*, itemsep=2pt]
  \item \texttt{C1}: Required roles present.
  \item \texttt{C2}: Type compatibility for each role.
  \item \texttt{C3}: Scope consistency of attachments/modifiers.
  \item \texttt{C4}: Tag precedence consistency with hard policy.
\end{itemize}

Typing alone is insufficient; well-formedness captures global coherence and policy-tag consistency.

\subsection{Policy Arbitration}
Surviving candidates are evaluated under:
\begin{itemize}[leftmargin=*, itemsep=2pt]
  \item \textbf{Soft rules}: defaults, conservative preferences, tie-breaks.
  \item \textbf{Hard rules}: non-negotiable safety/authority constraints.
\end{itemize}

Hard rules dominate soft score optimization.

\subsection{Confidence Scoring and Hard Gates}
For surviving candidates, composite confidence is:
\[
\mathrm{conf} = 0.20g + 0.25t + 0.25c + 0.20p + 0.10o
\]
where:
\begin{itemize}[leftmargin=*, itemsep=2pt]
  \item $g$: grammar adequacy
  \item $t$: type adequacy
  \item $c$: constraint satisfaction
  \item $p$: policy alignment
  \item $o$: override support
\end{itemize}

Hard gates are applied before final decision:
\begin{itemize}[leftmargin=*, itemsep=2pt]
  \item type failure $\rightarrow$ \texttt{REJECT}
  \item required-role failure $\rightarrow$ \texttt{REJECT}
  \item hard policy conflict $\rightarrow$ \texttt{REJECT}
\end{itemize}

Decision bands:
\begin{itemize}[leftmargin=*, itemsep=2pt]
  \item \texttt{ACCEPT} if $\mathrm{conf} \ge 0.78$
  \item \texttt{REVIEW} if $0.58 \le \mathrm{conf} < 0.78$
  \item \texttt{REJECT} if $\mathrm{conf} < 0.58$
\end{itemize}

The review band represents bounded automatic authority in ambiguous/policy-sensitive cases.

\subsection{Guarded Override and Conservative Fallback}
\begin{itemize}[leftmargin=*, itemsep=2pt]
  \item \textbf{Guarded override}: permits override only if override evidence itself passes typing, policy, and authorization checks.
  \item \textbf{Conservative fallback}: prefers safer admissible outcomes when ambiguity remains unresolved.
\end{itemize}

These mechanisms reduce unsafe overcommitment under uncertainty.

\subsection{Audit Trace Schema}
Each decision preserves:
\begin{itemize}[leftmargin=*, itemsep=2pt]
  \item input
  \item candidate set
  \item rules considered
  \item policy trace
  \item confidence breakdown
  \item score-selected candidate
  \item final-selected candidate
  \item reject reason (when applicable)
  \item override trace (when applicable)
\end{itemize}

This structure supports post hoc audit, dispute handling, and threshold calibration.

% ==================================================================
% 5. IMPLEMENTATION
% ==================================================================
\section{Implementation}

The runtime is implemented through ANIMA artifacts:
\begin{itemize}[leftmargin=*, itemsep=2pt]
  \item \texttt{ANIMA\_PARSER\_POLICY\_SPEC\_v0\_1\_plus.md}
  \item \texttt{parser\_policy\_example\_v0\_2.json}
  \item \texttt{ANIMA\_POLICY\_ENGINE\_EXTENSION\_v0\_3.md}
  \item \texttt{anima\_policy\_interpreter\_v0\_3.py}
  \item \texttt{calibration\_set\_v0\_3\_20.json}
  \item \texttt{run\_policy\_calibration\_v0\_3.py}
  \item \texttt{policy\_calibration\_v0\_3\_report.json}
\end{itemize}

Execution flow:
\begin{enumerate}[leftmargin=*, itemsep=2pt]
  \item input receive
  \item candidate proposal
  \item typed frame mapping
  \item constraint checks
  \item scoring
  \item score-selected determination
  \item final-selected arbitration
  \item action gate decision
  \item trace logging
\end{enumerate}

% ==================================================================
% 6. EVALUATION DESIGN
% ==================================================================
\section{Evaluation Design}

\subsection{Benchmark Families}
We evaluate four families:
\begin{enumerate}[leftmargin=*, itemsep=2pt]
  \item \textbf{Controlled ambiguity} (PP/RC/scope/role ambiguities)
  \item \textbf{Operational instruction} (urgency, authorization, check-skipping, bypass attempts)
  \item \textbf{Stability under paraphrase} (semantic equivalents with lexical/syntactic variation)
  \item \textbf{Auditability} (trace completeness and reconstructability)
\end{enumerate}

\subsection{Decision Buckets}
Calibration set buckets:
\begin{itemize}[leftmargin=*, itemsep=2pt]
  \item \texttt{clear\_accept}
  \item \texttt{clear\_reject}
  \item \texttt{review\_boundary}
  \item \texttt{override\_mixed}
\end{itemize}

\subsection{Metrics}
\begin{itemize}[leftmargin=*, itemsep=2pt]
  \item selection accuracy
  \item unsafe acceptance behavior
  \item hard reject correctness
  \item review-boundary handling
  \item decision stability
  \item audit completeness
\end{itemize}

% ==================================================================
% 7. RESULTS
% ==================================================================
\section{Results}

\subsection{Internal Calibration Results}
Using \texttt{policy\_calibration\_v0\_3\_report.json}:
\begin{itemize}[leftmargin=*, itemsep=2pt]
  \item total cases: 20
  \item passed cases: 20
  \item accuracy: 1.00
  \item decision distribution: \texttt{ACCEPT}~10, \texttt{REJECT}~6, \texttt{REVIEW}~4
\end{itemize}

\vspace{6pt}
\noindent Bucket accuracy:

\begin{table}[h]
\centering
\begin{tabular}{lcc}
\toprule
\textbf{Bucket} & \textbf{Result} & \textbf{Accuracy} \\
\midrule
\texttt{clear\_accept} & 8/8 & 1.00 \\
\texttt{clear\_reject} & 6/6 & 1.00 \\
\texttt{review\_boundary} & 4/4 & 1.00 \\
\texttt{override\_mixed} & 2/2 & 1.00 \\
\bottomrule
\end{tabular}
\caption{Calibration set bucket accuracy.}
\end{table}

\subsection{Interpretation}
The architecture shows internally consistent behavior across clear and boundary conditions, including override-sensitive cases. Hard-rule-first arbitration prevents inadmissible candidates from being promoted by confidence alone, and trace preservation enables post hoc reasoning reconstruction.

% ==================================================================
% 8. DISCUSSION
% ==================================================================
\section{Discussion}

\subsection{Plausibility Is Not Permissibility}
Generative quality remains valuable for proposal, but execution reliability requires post-generative admissibility controls.

\subsection{Review Is a Principled Outcome}
\texttt{REVIEW} is not failure; it is explicit bounded authority in uncertain or policy-sensitive regions.

\subsection{Runtime-Level Alignment}
Alignment is not only a model property; it is a runtime property. A well-tuned model in a weak runtime can still be operationally unsafe.

\subsection{Auditability as First-Class Reliability}
Explainable rejection and override trails are operational requirements, not optional diagnostics.

% ==================================================================
% 9. LIMITATIONS
% ==================================================================
\section{Limitations}

\begin{enumerate}[leftmargin=*, itemsep=2pt]
  \item Reported results are internal calibration results, not full external generalization evidence.
  \item Some policy rules and thresholds are deployment-domain dependent.
  \item Candidate proposal quality remains an upstream dependency.
  \item Human-factors usability of traces requires dedicated reviewer studies.
\end{enumerate}

% ==================================================================
% 10. CONCLUSION
% ==================================================================
\section{Conclusion}

This paper introduced a layered neuro-symbolic runtime that formalizes the boundary between plausible interpretation and permissible execution. By combining typed semantic filtering, well-formedness constraints, policy gating, guarded override, conservative fallback, and explicit audit tracing, the runtime converts candidate interpretations into semantically valid, policy-aligned, and inspectable operational decisions.

The main takeaway is straightforward:

\begin{quote}
\emph{What a model can plausibly say is not the same as what a system should permissibly do.}
\end{quote}

For action-coupled AI systems, post-generative control layers are essential infrastructure.

% ==================================================================
% METHODS APPENDIX
% ==================================================================
\appendix
\section{Methods Appendix (Condensed)}

\subsection{Candidate Representation}
Each candidate includes parse/interpretation, semantic frame, type result, constraints result, policy score, override status, and final decision.

\subsection{Thresholding}
Hard constraints first, then confidence banding into \texttt{ACCEPT}/\texttt{REVIEW}/\texttt{REJECT}.

\subsection{Comparative Protocol Note}
This version reports internal calibration evidence only. Claims about superiority over external baselines require additional benchmarked experiments and tabled results.

% ==================================================================
% DATA AND CODE AVAILABILITY
% ==================================================================
\section*{Data and Code Availability}

Current artifact paths are maintained in the ANIMA workspace and can be packaged for review according to repository policy.

% ==================================================================
% REFERENCES
% ==================================================================
\section*{References}

Representative references include core work on transformers, semantic parsing, causal reasoning, explainability, neuro-symbolic methods, and AI safety/governance. References will be finalized with proper citation-format bibliography prior to submission.

\end{document}
